\section{The cubic trial and its exact residual}\label{sec:trial}
On each seed torus let $D$ denote its actual adjacency, including wraparound
edges. Put $D_2=\off(D^2)=D^2-4I$ and define
\begin{equation}\label{eq:K}
 K=-\frac g4D+\frac{g^2}{8}(D_2-D)
       -\frac{g^3}{64}(5D^3-4D_2-40D).
\end{equation}
This symmetric matrix has zero diagonal. Define
\begin{equation}\label{eq:trial}
 S_K=\frac12\sum_{r,s,\alpha}K_{rs}x_{r\alpha}x_{s\alpha},
 \qquad \psi_K=e^{S_K}\vac
   =\sum_{j=0}^{\lfloor A/2\rfloor}\frac{S_K^j}{j!}\vac.
\end{equation}
The last equality is exact hard-core nilpotence. The vacuum coefficient is one,
so $\norm{\psi_K}\ge1$. This is a specified polynomial trial state at every
$g$, not a discarded-tail perturbation expansion.

For one edge $\{r,s\}$ put $x_\alpha=x_{r\alpha}$, $y_\alpha=x_{s\alpha}$,
$k=K_{rs}$ and
$u_\alpha=\sum_{t\ne r,s}K_{rt}x_{t\alpha}$,
$v_\alpha=\sum_{t\ne r,s}K_{st}x_{t\alpha}$. All these variables commute;
write $u\cdot v=\sum_\alpha u_\alpha v_\alpha$. Define the following
\emph{literal creation polynomials}:
\begin{align}\label{eq:q2}
 q_2={}&u\cdot v+\sum_\alpha x_\alpha(v_\alpha-3ku_\alpha)
       +\sum_\alpha y_\alpha(u_\alpha-3kv_\alpha)
       +(1-2k-3k^2)\sum_\alpha x_\alpha y_\alpha,\\
 q_4={}&-\sum_\alpha x_\alpha u_\alpha(u\cdot v)
         -\sum_\alpha y_\alpha v_\alpha(u\cdot v)\notag\\
 &+\sum_{\alpha,\gamma}x_\alpha y_\gamma
 \left[v_\alpha u_\gamma-\delta_{\alpha\gamma}u\cdot v
       -u_\alpha u_\gamma-v_\alpha v_\gamma
       +3ku_\alpha v_\gamma-k\delta_{\alpha\gamma}u\cdot v\right],\label{eq:q4}\\
 q_6={}&\sum_{\alpha,\gamma}x_\alpha y_\gamma
                u_\alpha v_\gamma(u\cdot v).\label{eq:q6}
\end{align}
Each has its indicated creation degree; powers of $g$ in its coefficients are
not being truncated.

\begin{proposition}[Exact transformed Hamiltonian]\label{prop:residual}
For any symmetric zero-diagonal $K$, define
\[
 e_A=\frac{3g}{2}\sum_{\{r,s\}\in\mathcal E}K_{rs},\quad
 R_2=2S_K+\frac g2\sum_{\mathcal E}q_2,\quad
 R_4=\frac g2\sum_{\mathcal E}q_4,\quad
 R_6=\frac g2\sum_{\mathcal E}q_6.
\]
Then
\begin{equation}\label{eq:residual}
 (H^{\rm C}-e_A)\psi_K
       =e^{S_K}(R_2+R_4+R_6)\vac
       =(R_2+R_4+R_6)\psi_K.
\end{equation}
For \eqref{eq:K}, the three reference shifts are
\begin{align}\label{eq:shifts}
 e_{16}&=-12g^2-6g^3-\tfrac{15}{2}g^4,\notag\\
 e_{32}&=-24g^2-12g^3-\tfrac{45}{4}g^4,\\
 e_{64}&=-48g^2-24g^3-15g^4.\notag
\end{align}
\end{proposition}
\begin{proof}
Split $S_K$ into its outside-site part and the two-endpoint polynomial
$\sum_\alpha(x_\alpha u_\alpha+y_\alpha v_\alpha+kx_\alpha y_\alpha)$.
The outside part commutes with the bond. The exact local dressed vacuum has
components $1$, $u_\alpha$, $v_\gamma$, and
$u_\alpha v_\gamma+k\delta_{\alpha\gamma}$.
Apply \eqref{eq:bondaction}, then the inverse dressing obtained by negating
$u,v,k$. Collecting the sixteen resulting entries gives
$3k+q_2+q_4+q_6$ for twice the transformed bond. Also
$e^{-S_K}Ne^{S_K}\vac=2S_K\vac$.
Summation proves the first identity. Every $R_d$ is in the same commutative
creation algebra as $S_K$, proving the second. Evaluating the edge entries of
the explicitly given polynomial in $D$ yields \eqref{eq:shifts}; the different
fourth-order coefficients retain the actual short wrapping walks.
\end{proof}

The reference $e_A$ is \emph{not} the Rayleigh quotient. In particular, the
nonzero degree-six term in \eqref{eq:q6} is necessary. For the cubic kernel the
coefficient degrees reach $g^7$, $g^{10}$, and $g^{13}$ in $R_2,R_4,R_6$,
respectively, but the maximum creation degree remains six.
